% Transcription — fonds Alexandre Grothendieck, shelfmark 68, batch 2 (pages 21-40).
% Established from the facsimile archives/batches/68/batch-02.pdf (a mirror of
% https://grothendieck.umontpellier.fr/68.pdf), per the skill
% transcribe-grothendieck. The batch file carries no cover sheet: PDF page N is
% archivists' page N + 20.
%
% Pass: Opus 5.5 (claude-opus-5-5), 2026-09-22 — first pass, unchecked against the pages by a human.
%
% What the batch is.
%   21     End of a run begun in batch 1 (pp. 19-20 there: « 3 Couples
%          prétassés », « 4. Couples tassés »): positive functions of finite
%          support on a group G, convolution *, the relations f' ≺ f and
%          f' ~ f, level sets E_α(f), P_f(G). Conditions d), e) on pairs and
%          triples of finite subsets in S, then a set S stable by translation
%          with S_0 = S ∩ P_f(G) and a decomposition f = Σ α_i A_i over a
%          totally ordered family. Blue ink. An isolated 0 → _2G → G → G at
%          the foot.
%   22     Blank. Absent.
%   23-33  « Jeux de position » (his heading, p. 23): positional games —
%          configurations C, move relation R ⊂ C × C, terminal set C_0,
%          players J, initiative α : C ∖ C_0 → J, winning terminal sets
%          G_0(j); parties, strategies, winning positions G(j); the
%          fixed-point theorem for G(j) (p. 25) proved by a well-ordering
%          argument (pp. 26-28); strategies « contre j » and admissible
%          positions (pp. 28-29); for two players, the partition
%          C = G(j) ⊔ G(j') ⊔ N; alternating games, α̂ : C → J, the sets
%          G, P, N of winning, losing and null positions (pp. 30-32); a
%          characterisation of (G, P) when there is no infinite partie
%          (p. 33, blue ink). Pencil, fast hand, pp. 24-29 faint.
%   34     Blank. Absent.
%   35-40  « Exponentielles (en car. 0) »: I) exp / log on the nilradical of
%          a Q-algebra; II) one-parameter formal groups: homomorphisms
%          Ĝ_a → Un(U) and G_a → Un(U) are e_x, x ∈ U resp. x ∈ N(U)
%          (Th. 2); the same for Aut_K(M); morphisms of G_a-modules
%          (pp. 39-40). Page 40 breaks off mid-proof (« implique par
%          récurrence »); the argument continues into batch 3.
%
% Notation fixed for the batch: R for his relation (drawn as a script R on
% p. 23, plain afterwards); \mathfrak{P}(C) for the power set he draws as a
% gothic P; \overline{C}(j) for his C barre; \complement for his ∁;
% \hat\alpha and \bar{\hat\alpha}; \mathbb{A} for the outlined A of p. 35;
% \mathcal{U} for his U; \overline{\mathbb{G}}_{a,K} for the formal line he
% writes with a bar; \mathbf{Aut}, \mathbf{End} for the functors he draws
% with outlined initials.
%
% The dating is the inventory's own for the folder, copied verbatim. Its
% group, [66-89] « Autour de l'enseignement - Thèmes liés (cours et
% directions de recherches) », is dated 1964-1984.
\documentclass[11pt,a4paper]{article}
% Le préambule commun à toutes les transcriptions du fonds Grothendieck.
%
% Il tient en une idée : **l'incertitude se compose, elle ne se résout pas.**
% Une transcription qui lisse un mot douteux a détruit la seule chose pour
% laquelle on la faisait — un lecteur doit pouvoir savoir, à chaque endroit,
% ce qui était sur le feuillet et ce qui a été ajouté. Les six macros qui
% suivent sont donc l'appareil critique, et non de la décoration.
%
% Les mêmes commandes sont reconnues par `scripts/render.mjs`, qui produit la
% vue de lecture du site. Une macro ajoutée ici doit l'être aussi là-bas,
% faute de quoi le rendu échouera bruyamment — ce qui est voulu : un rendu
% silencieusement faux est pire qu'un rendu absent.

\ProvidesPackage{grothendieck}[2026/08/08 Transcriptions du fonds Grothendieck]

\usepackage{amsmath,amssymb,amsthm}
\usepackage{mathrsfs}
\usepackage{stmaryrd}
% Les diagrammes commutatifs sont partout dans ce fonds : c'est la langue
% même de la géométrie algébrique de Grothendieck, et une transcription qui
% les rendrait en prose ne transcrirait rien.
\usepackage{tikz-cd}
% Français par défaut — c'est la langue des feuillets — mais l'anglais est
% chargé aussi : la traduction et le résumé sont des documents anglais, et la
% typographie française (espaces avant les deux-points) y serait une faute.
% Ces documents basculent par \selectlanguage{english} en tête de corps.
\usepackage[english,french]{babel}
\usepackage{fontspec}
\usepackage{geometry}
\geometry{a4paper,margin=25mm,marginparwidth=28mm}
\usepackage{marginnote}
\usepackage{xcolor}
% ulem, not soul. `soul`'s \st and \ul are built on a character-by-character
% scan that XeLaTeX's font machinery breaks: the document fails with an
% undefined control sequence rather than degrading. ulem does the same two jobs
% and is Unicode-engine safe. `normalem` stops it hijacking \emph, which the
% transcriptions use for Grothendieck's own emphasis.
\usepackage[normalem]{ulem}
\usepackage{hyperref}
\hypersetup{colorlinks=true,linkcolor=black,urlcolor=black,pdfborder={0 0 0}}

% --- Métadonnées du lot -----------------------------------------------------
% Reprises telles quelles par le script de rendu, qui les affiche en tête de
% la vue de lecture. Elles fixent ce que le fichier prétend couvrir : sans
% elles, une transcription flotte sans point d'attache dans un fonds de seize
% mille feuillets.

\newcommand{\folder}[1]{\gdef\grothFolder{#1}}
\newcommand{\batch}[1]{\gdef\grothBatch{#1}}
\newcommand{\leaves}[2]{\gdef\grothFirst{#1}\gdef\grothLast{#2}}
\newcommand{\foldertitle}[1]{\gdef\grothFoldertitle{#1}}
\gdef\grothFolder{?}\gdef\grothBatch{?}\gdef\grothFirst{?}\gdef\grothLast{?}\gdef\grothFoldertitle{}

% --- Le résumé --------------------------------------------------------------
%
% Il ouvre la lecture modernisée, sous le titre « Résumé », et il est composé
% en petit. Ce n'est pas un ornement : c'est le seul passage du document qui
% ne soit pas des mathématiques, et un lecteur doit pouvoir le reconnaître
% avant de l'avoir lu — puis le sauter s'il connaît déjà le sujet, ou s'y
% arrêter s'il ne le connaît pas. Le filet qui le suit marque la reprise du
% corps.

\newenvironment{resume}
  {\small\setlength{\parskip}{0.45em}}
  {\par\medskip\hrule\medskip}

% --- Mots-clés ---------------------------------------------------------------
%
% En anglais, en clôture du résumé de la lecture modernisée : le vocabulaire
% moderne sous lequel chercher ce que les feuillets construisent. Le manifeste
% du site (`npm run manifest`) les extrait pour étiqueter la cote entière —
% cette ligne est la source unique des tags, il n'y a pas de fichier à côté.
\newcommand{\keywords}[1]{\par\smallskip\noindent
  {\footnotesize\textsc{Keywords} --- \textit{#1}}\par}

% --- L'appareil critique ----------------------------------------------------

% Le numéro de feuillet, en marge. C'est l'ancre qui permet de régler un
% désaccord sur une formule en regardant *un* feuillet plutôt qu'une chemise
% de six cents. Le site s'en sert aussi pour tourner les pages du fac-similé
% quand on fait défiler la transcription.
\newcommand{\leaf}[1]{%
  \marginnote{\footnotesize\color{gray!70}\textsf{#1}}%
  \label{leaf:#1}\ignorespaces}

% Illisible : on ne devine pas. Un mot inventé qui se lit comme les autres est
% le pire résultat possible.
\newcommand{\ill}{\textcolor{red!60!black}{[\,\ldots\,]}}

% Lecture douteuse : le mot est donné, et signalé comme incertain.
\newcommand{\uncertain}[1]{\uline{#1}}

% Ajout de l'éditeur — une lettre manquante, un mot sous-entendu.
\newcommand{\add}[1]{\textcolor{blue!50!black}{[#1]}}

% Biffé par Grothendieck lui-même. Ce qu'il a rayé fait partie du document :
% une rature dit souvent où la pensée a changé de direction.
\newcommand{\struck}[1]{\sout{#1}}

% Note du transcripteur, sur l'état matériel du feuillet.
\newcommand{\note}[1]{\par\smallskip{\small\color{gray!60!black}\textit{[#1]}}\par\smallskip}

% Note marginale de Grothendieck, distincte du corps du texte.
\newcommand{\marginal}[1]{\par\smallskip\noindent\hspace{1em}%
  {\small\color{gray!60!black}#1}\par\smallskip}

% --- Titre ------------------------------------------------------------------

\AtBeginDocument{%
  \begin{center}
    {\large\bfseries Fonds Alexandre Grothendieck --- cote n\textsuperscript{o}~\grothFolder}\\[2pt]
    {\itshape\grothFoldertitle}\\[4pt]
    {\small feuillets \grothFirst--\grothLast\ (lot \grothBatch)}\\[2pt]
    {\footnotesize\color{gray!60!black}Transcription établie d'après le fac-similé numérisé
     par l'Université de Montpellier.\\
     Les lectures incertaines sont soulignées, les illisibles notées [\,\ldots\,],
     les ajouts entre crochets.}
  \end{center}
  \vspace{1em}}

\endinput


\folder{68}
\batch{2}
\pages{21}{40}
\dating{[à partir de 1978-à partir de 1983]}
\watermark{Édition de démonstration}
\foldertitle{Jeux de position : notes manuscrites (s.d.)}

\begin{document}

\page{21}
\note{la page continue un développement commencé au lot précédent (pp.~19-20,
« Couples prétassés », « Couples tassés ») : $\mathbb{R}[G]^{+}$, produit de
convolution $*$, relations $f' \prec f$ et $f' \sim f$, ensembles
$E_\alpha(f)$. Les conditions a) à c) sont à la page 20.}

\marginal{si (K) \uncertain{les conditions} d) et si les \uncertain{couples}
$\in S \times S$ sont prétassés $\downarrow$ \uncertain{id.}}

d) $\forall A, B \in \mathfrak{P}_f(G) \cap S$, $A' \sim A$, $B' \sim B$, on a
\[
A' * B' \prec A * B
\]
i.e.\ \emph{$(A, B)$ tassé}.

e) $\forall A, B, C \in \mathfrak{P}_f(G) \cap S$, $A' \sim A$, $B' \sim B$,
$C' \sim C$, on a
\[
\| A' * B' * C' \|_\infty \leq \| A * B * C \|_\infty
\]

\struck{Supposons}
Conclu\uncertain{t} : Si l'on a $S$, \uncertain{avec} $S$ stable par
translation, $S_0 = S \cap \mathfrak{P}_f(G)$, \uncertain{avec}
$S_0 \neq \emptyset$. Soit $f \in S$. Alors
\begin{itemize}
  \item[a)] $\exists\, t \in G$, $f = \varepsilon_t * f$
  \note{tel qu'écrit ; l'accent sur l'un des deux $f$ n'est pas visible. Le
  crochet qui suit et la p.~19 (« $\check f = \varepsilon_s * f$ ») font
  attendre $\check f = \varepsilon_t * f$.}
  \quad [$(f, f)$ prétassé]
  \item[b)] \ill{} $E_\alpha(f) \in S_0$ \quad [$(f, A)$ prétassé pour
  $A \in S_0$]
\end{itemize}

\struck{Inversement, soit $f \in \mathbb{R}[G]^{+}$}
Soit \struck{$S$}
\[
\widetilde{S} = \{ f \in \mathbb{R}(G)^{+} \mid \text{a), \ill{} b)} \}
\]
\struck{$f \in \widetilde S$} i.e.
\[
\begin{cases}
f = \sum_I \alpha_i A_i, & \alpha_i > 0,\quad A_i \in S_0 \quad (A_i) \text{ tot.\ ord.} \\
\exists\, s \in G, & A_i = \varepsilon_s * \check{A}_i \quad \forall i
\end{cases}
\]
Soit de même
\[
\begin{cases}
g = \sum_J \beta_j B_j, & \beta_j > 0,\quad B_j \in S_0 \quad (B_j) \text{ tot.\ ord.} \\
t \in G, & B_j = \varepsilon_t * \check{B}_j
\end{cases}
\]
On a
\[
f * g = \sum_{(i,j) \in I \times J} \alpha_i \beta_j \, A_i * B_j
\]
\struck{Supposons que $\exists\, (G_n)$}
On a $A_i * B_j \in S$, \uncertain{donc} \ill{}
\[
E_*(A_i * B_j) \subset S_0
\]
\textbf{Supposons} $\bigcup E_\alpha(A_i * B_j)$ \emph{tot.} ordonné ; on
\uncertain{veut}
\[
\bigcap_{C \in \bigcup E_\alpha(A_i * B_j)} C \neq \emptyset
\]
\note{isolé au pied de la page, sans lien apparent avec ce qui précède :}
\[
0 \to {}_2 G \to G \xrightarrow{\;2\;} G
\]

\section{Jeux de position}
\note{le titre est de sa main, en tête de la page 23.}

\page{23}
a) $C$ ensemble des “\emph{configurations}” ou “\emph{positions}”

b) $R \subset C \times C$, interprété comme $x \mapsto R(x)$,
$C \to \mathfrak{P}(C)$,
\note{il écrit ici un $R$ rond ; les pages suivantes ont un $R$ droit, que
l'on garde partout.}
d'où
\[
C_0 = \{ x \in C \mid R(x) = \emptyset \} \subset C \qquad (\text{conf.\ \textit{terminales}})
\]

c) $J$ ensemble des “\emph{joueurs}”

d) $C \setminus C_0 \xrightarrow{\;\alpha\;} J$. On dit que (si
$x \in C \setminus C_0$) $j = \alpha(x)$ est le joueur qui a
\emph{l'initiative} dans la position $x$.
\[
\left.
\begin{array}{l}
C(j) = \{ x \in C \setminus C_0 \mid \alpha(x) = j \} \\
\overline{C}(j) = (C \setminus C_0) \setminus C(j)
\end{array}
\right\}
\]

e) $J \to \mathfrak{P}(C_0)$, $j \mapsto G_0(j)$ (ens.\ des \emph{positions
terminales gagnantes pour $j$})

\marginal{ne dépend que de $C$, $R$}
\emph{Partie} $(x_0, x_1, \dots, x_n, \dots)$ suite finie ou infinie de
positions, telles que pour deux termes consécutifs $x_i, x_{i+1}$, on ait
\[
x_{i+1} \in R(x_i)
\]
\emph{Prolongement} des parties — parties non prolongeables ou
“\emph{terminées}” : ce sont les parties qui sont soit infinies, soit finies,
avec \struck{comme} \add{la} position finale $x_n \in C_0$.
\marginal{Composition des parties (Associatif et \ill{} — \ill{} aux
“catégorie des positions” …)}

\marginal{ne dépend que de $C$, $R$ et $G_0(j) \subset C_0$}
\emph{Parties gagnées pour $j$} : partie finie (non prolongeable) dont la
position \struck{finale} finale $x_n$ est $\in G_0(j)$.

\marginal{NB Si $XY$ est gagnée par $j$, $\Leftrightarrow$ $Y$ est gagnée
par $j$.}

\emph{Stratégie $\Sigma$ pour $j$} (ou $j$-stratégie) : (Application
$\Sigma : C(j) \to \mathfrak{P}(C)$ telle que $\forall x \in C(j)$, on ait

\page{24}
\[
\Sigma(x) \subset R(x), \qquad \Sigma(x) \neq \emptyset .
\]
Cette notion ne dépend que de $C$, $R$, $C(j) \subset C \setminus C_0$.

\emph{NB} On pose $\Sigma(x) = R(x)$ si $x \in \overline{C}(j)$. Donc on peut
\struck{\ill{}} \struck{interpréter $\Sigma$ comme} une application
$C \setminus C_0 \to \mathfrak{P}^{*}(C)$, avec $\Sigma(x) \subset R(x)$ et
$\Sigma(x) \neq \emptyset$ $\forall x \in C \setminus C_0$, et
$\Sigma(x) = R(x)$ si $x \in \overline{C}(j)$.

Partie jouée en accord avec une stratégie $\Sigma$ (ou
\emph{compatible} avec $\Sigma$, ou $\Sigma$-partie) : cette notion est stable
par composition. (On a une sous-catégorie \uncertain{comp.\ à} ci-dessus,
mais a priori moins de flèches)

\emph{Lemme} Soit \struck{$\Sigma$} $J' \subset J$, $\forall j \in J'$ soit
$\Sigma_j$ une $j$-stratégie. Soit $x_0 \in C$. Alors $\exists$ partie
maximale d'origine $x_0$, compatible avec les $\Sigma_j$ ($j \in J'$).

\marginal{ne dépend que de $C$, $R$, $C(j)$, $G_0(j)$}
Stratégie $\Sigma$ (pour $j$) \emph{gagnante pour $x_0 \in C$} (ou
\emph{position $x_0$ gagnante} pour une $j$-stratégie $\Sigma$) : toute partie
maximale commençant en $x_0$, et jouée suivant $\Sigma$, est gagnée par $j$
(i.e.\ est finie, et son extrémité $x_n \in C_0$ est $\in G_0(j)$).

\emph{Position} $x_0 \in C$ \emph{gagnante pour $j$} : telle qu'il existe une
$j$-stratégie qui soit gagnante pour $x_0$. Soit
\[
G(j)
\]
l'ens.\ de ces positions. On a
\[
G(j) \cap C_0 = G_0(j)
\]

[\emph{Proposition} Soit $J' \subset J$. Si
$\bigcap_{j \in J'} G(j) \neq \emptyset$, alors
$\bigcap_{j \in J'} G_0(j) \neq \emptyset$. En d'autres termes, si
$\bigcap_{j \in J'} G_0(j) = \emptyset$, alors
$\bigcap_{j \in J'} G(j) = \emptyset$ i.e.\ une $x_0 \in C$ ne peut être
gagnante pour tous les $j \in J'$.

Immédiat par le Lemme.]

\page{25}
\emph{Proposition} Soient $\Sigma$ une $j$-stratégie, $x_0 \in C$
\struck{supposons} gagnante pour $\Sigma$,
$X = (x_0, x_1, \dots, x_n, \dots)$ une partie (d'origine $x_0$) jouée suivant
$\Sigma$, alors toutes les positions $x_i$ sont gagnantes pour $\Sigma$.

Soit en effet $Y = (y_0 = x_i, y_1, \dots, y_m, \dots)$ une partie
\uncertain{maximale} jouée suivant $\Sigma$, alors $XY$ est jouée suivant
$\Sigma$ et est d'origine $x_0$ gagnante pour $\Sigma$, donc est
\struck{gagnée par partie finie, telle $Y$} \struck{a finie} donc $Y$ est gagnée
par $j$.
\note{les deux ratures de cette phrase se chevauchent ; la lecture des mots
barrés est incertaine.}

\emph{Th} On a \struck{\ill{}}
\[
\begin{array}{ll}
\text{a)} & G(j) \cap C(j) = \{ x_0 \in C(j) \mid R(x_0) \cap G(j) \neq \emptyset \} \\
\text{b)} & G(j) \cap \overline{C}(j) = \{ x_0 \in \overline{C}(j) \mid R(x_0) \subset G(j) \}
\end{array}
\]
[et pour mémoire
\[
G(j) \cap C_0 = G_0(j) \quad ].
\]

\emph{Dém} a) On suppose $x_0 \in C(j)$.

1) $\subset$ Soit $x_0 \in G(j)$ \struck{\ill{}}, soit $\Sigma$ une stratégie
pour $j$ gagnante pour $x_0$, considérons $x_1 \in \Sigma(x_0)$, alors
$(x_0, x_1)$ est une $\Sigma$-partie, donc par la prop.\ précédente, $x_1$ est
gagnante pour $\Sigma$, donc gagnante, donc $x_1 \in R(x_0) \cap G(j)$, donc
$R(x_0) \cap G(j) \neq \emptyset$.

2) $\supset$ Soit $x_1 \in R(x_0) \cap G(j)$. Prouvons que $x_0 \in G(j)$.
\struck{Considérons} Soit $\Sigma_1$ une $j$-stratégie $x_1$-gagnante. Si
$x_0$ est gagnante pour $\Sigma_1$, on a $x_0 \in G(j)$, OK. Sinon, définissons
une $j$-stratégie $\Sigma_0$ $x_0$-gagnante ainsi :
\[
\begin{array}{l}
\Sigma_0(x_0) = \{ x_1 \} \\
\Sigma_0(x) = \Sigma_1(x) \quad \text{si } x \neq x_0
\end{array}
\]

\page{26}
C'est une $j$-stratégie, \struck{prouvons} Prouvons qu'elle est gagnante pour
$x_0$. Une partie maximale jouée suivant $\Sigma_0$ \struck{commençant} avec
$(x_0, x_1)$, soit $(x_0, x_1, \dots, x_n, \dots)$ \struck{les $x_i$
($i \geq 1$) sont gagnants pour $\Sigma_1$} prouvons qu'elle est gagnée par
$j$. Prouvons \struck{que tous} que les $x_i$ ($i \geq 1$) sont $\neq x_0$.
Sinon, \struck{si $i$ est le} soit le plus petit indice $i \geq 1$ tel que
$x_i = x_0$, \struck{pour un $i \geq 1$}, alors \struck{par la prop.\
impliquerait} [$(x_1, \dots, x_i = x_0)$ serait compatible avec $\Sigma_1$],
on trouverait que $x_i = x_0$ est gagnant pour $\Sigma_1$, contrairement à
l'hypothèse que $x_0$ non gagnant pour $\Sigma_1$.

Comme les $x_i$ ($i \geq 1$) sont $\neq x_0$, la partie $(x_1, x_2, \dots)$
est une $\Sigma_1$-partie, donc elle est gagnée par $j$, donc
$(x_0, x_1, \dots)$ aussi, cqfd.

b) On suppose $x_0 \in \overline{C}(j)$.

1) $\subset$. Supposons $x_0 \in G(j)$, prouvons $R(x_0) \subset G(j)$. Soit
$x_1 \in R(x_0)$, prouvons $x_1 \in G(j)$. En effet, [si $\Sigma_0$ est une
$j$-stratégie $x_0$-gagnante,] $(x_0, x_1)$ est une $\Sigma_0$-partie, donc
par la prop., $x_1$ est gagnante pour $\Sigma_0$, donc $x_1 \in G(j)$.

2) $\supset$. Supposons $R(x_0) \subset G(j)$, prouvons $x_0 \in G(j)$.
$\forall x_1 \in R(x_0)$, soit $\Sigma_{x_1}$ une $j$-stratégie
$x_1$-gagnante. Choisissons un \emph{bon ordre} sur $R(x_0)$. On va définir
une $j$-stratégie $\Sigma$ $x_0$-gagnante ainsi.

$\forall x \in C$, soit
\[
E(x) = \{ x_1 \in R(x_0) \mid \exists\ \Sigma_{x_1}\text{-partie de } x_1 \text{ à } x \}
\]
\note{l'indice du point de départ, dans « de $x_1$ à $x$ », est surchargé ;
$x_1$ est la lecture que demande la suite.}

\page{27}
Soit
\[
C' = \{ x \in C \mid E(x) \neq \emptyset \}
\]
et pour $x \in C'$, posons
\[
\varepsilon(x) = \text{plus petit élément de } E(x)
\]
(donc $\varepsilon : C' \to R(x_0)$). Posons
\[
\Sigma_0(x) =
\begin{cases}
R(x) & \text{si } x \in \overline{C}(j) \cup \complement C' \\
\Sigma_{\varepsilon(x)}(x) & \text{si } x \in C(j) \cap C'
\end{cases}
\]
Donc $\Sigma_0$ est une $j$-stratégie. On va prouver qu'elle est gagnante pour
$x_0$.

\emph{Lemme} \struck{Soit $x \in C'$} Si $x \in C'$, $x$ est gagnant pour
$\Sigma_{\varepsilon(x)}$ (donc gagnant pour $j$ — donc $C' \subset G(j)$) et
\struck{$C' \subset G(j)$, et si $x \in C'$, $y \in \Sigma_0(x)$, alors
$\Sigma_0(x) \subset C'$ et} pour tout $y \in \Sigma_0(x)$, on a
\[
E(y) \ni \varepsilon(x)
\]
(donc $E(y) \neq \emptyset$, donc $y \in C'$, donc $\Sigma_0(x) \subset C'$
$[\subset G(j)]$) et par suite
\[
\varepsilon(y) \leq \varepsilon(x).
\]

\emph{Dém.} Soit $x \in C'$ donc $E(x) \neq \emptyset$, alors si
$x_1 \in E(x)$, $\exists$ \struck{partie} $\Sigma_{x_1}$-partie de $x_1$ à $x$,
donc par la prop.\ (comme $\Sigma_{x_1}$ est $x_1$-gagnante), $x$ est
$\Sigma_{x_1}$-gagnante, donc a fortiori $x \in G(j)$.

Soit $y \in \Sigma_0(x) = \Sigma_{\varepsilon(x)}(x)$ ;
$(x_1, x_2, \dots, x_n = x, x_{n+1} = y)$, où $x_1 = \varepsilon(x)$, est une
$\Sigma_{\varepsilon(x)}$-partie, donc \struck{\ill{}}
$\varepsilon(x) \in E(y)$, OK.

Prouvons que $\Sigma_0$ est gagnante pour $x_0$. Soit une partie maximale
commençant par $x_0$

\page{28}
et comp.\ avec $\Sigma_0$
\[
(x_0, x_1, x_2, \dots)
\]
Par le lemme, il vient de proche en proche \struck{cette partie est}
(puisque $x_1 \in C'$) que \struck{$x_1 \in G(j)$,} les $x_i$ ($i \geq 1$)
sont $\in C'$ \struck{\ill{}}, donc $j$-gagnants. Pour montrer que la partie
est gagnée par $j$, il suffit de prouver qu'elle est finie. \struck{Or}
Supposons le contraire. Les $\varepsilon(x_i)$ ($i \geq 1$) forment une suite
décroissante d'éléments de $R(x_0)$, donc stationnaire. Soit $n$ tel que
$\varepsilon(x_m) = \varepsilon(x_n)$ $(= \xi)$ pour $m \geq n$. Alors par
définition de $\Sigma_0$, la partie $(x_n, x_{n+1}, \dots)$ est une partie
comp.\ avec $\Sigma_\xi$. Comme $x_n$ est gagnant pour $\Sigma_\xi$ (par le
lemme), cette partie est finie, absurde.

Stratégie \emph{contre $j$} (stratégie relative à : $\overline{C}(j)$), i.e.\
$\Sigma'$ telle que $\Sigma'(x) = R(x)$ si $x \in C(j)$ (et
$\Sigma'(x) \subset R(x)$, $\Sigma'(x) \neq \emptyset$ si
$x \in \overline{C}(j)$). Admissible pour $x_0$ (toute partie maximale
d'origine $x_0$ jouée suivant $x_0$
\note{sic ; on attend « suivant $\Sigma'$ ».}
est \emph{non gagnée} par $j$, i.e.\ est ou bien infinie, ou bien a une
position terminale $\notin G_0(j)$).

\struck{Corollaire} \emph{Théorème} Soit $\Sigma'$ la stratégie contre $j$
définie par

\page{29}
\[
\Sigma'(x) =
\begin{cases}
R(x) & \text{si } x \in C(j) \cup G(j) \\
R(x) \cap \complement G(j) & \text{si } x \in \overline{C}(j) \cap \complement G(j)
\end{cases}
\]
(C'est une stratégie contre $j$, par [b) $\supset$]). Alors pour $x \in C$,
les conditions suivantes sont équivalentes
\begin{itemize}
  \item[a)] $x_0 \notin G(j)$
  \item[b)] $x_0$ est admissible contre $j$
  \item[c)] $x_0$ est admissible pour $\Sigma'$ contre $j$.
\end{itemize}

\emph{Dém} c) $\Rightarrow$ b) trivial, b) $\Rightarrow$ a) par le Lemme.
Prouvons a) $\Rightarrow$ c). Si \ill{} une partie $(x_0, x_1, \dots)$
d'origine $x_0$ jouée en accord avec $\Sigma'$, alors par définition de
$\Sigma'$, on voit de proche en proche que $x_i \in \complement G(j)$
$\forall i$, donc la partie ne peut être gagnée par $j$ !

\emph{Corollaire} Supposons $J = \{ j, j' \}$, $j \neq j'$, et
\[
G_0(j) \cap G_0(j') = \emptyset .
\]
Soit \struck{\ill{}}
\[
N = \{ x \in C \mid x \text{ admissible contre } j, \text{ et contre } j' \}.
\]
Alors $C$ est réunion disjointe de $G(j)$, $G(j')$ et $N$.

\page{30}
\struck{Supposons}
\emph{NB} $G(j)$, $G(j')$ et $N(j)$ sont connus quand on connaît (en plus des
parties $G_0(j)$ et $G_0(j')$ de $C_0$) les ensembles
\[
\begin{cases}
G^{*} = \{ x \in C \setminus C_0 \mid x \in G(\alpha(x)) \} \\
P^{*} = \{ x \in C \setminus C_0 \mid x \in G(\bar\alpha(x)) \}
\end{cases}
\]
où $\bar\alpha(x)$ désigne l'unique $j' \in J$ tel que
$j' \neq j = \alpha(x)$.

On dit qu'une position $x \in C \setminus C_0$ qui est $\in G^{*}$ (resp.\
$\in P^{*}$), i.e.\ \struck{telles} gagnante pour le joueur qui a
l'initiative (resp.\ pour le joueur qui n'a pas l'initiative) est
\emph{gagnante} (resp.\ \emph{perdante}) ; on dit que c'est une
\emph{position nulle} \struck{autrement} si elle n'est ni gagnante, ni
perdante, et on pose
\[
N^{*} = (C \setminus C_0) \setminus (G^{*} \cup P^{*}) = N \cap (C \setminus C_0)
\]
(où $N$ est défini dans le corollaire ci-dessus).

On aura
\[
\begin{cases}
G(j) \cap C(j) = G^{*} \cap C(j) \\
G(j) \cap \overline{C}(j) = G(j) \cap C(j') = P^{*} \cap C(j') \\
G(j) \cap C_0 = G_0(j)
\end{cases}
\]
ce qui montre que les $G(j)$ ($j \in J$) se reconstituent par la connaissance
de $G^{*}$, $P^{*}$ et des $G_0(j)$ ($j \in J$).

Supposons qu'on ait une application
\[
\hat\alpha : C \to J
\]
prolongeant $\alpha$, et qu'on ait, pour tout $x \in C$ et $y \in R(x)$
\[
\hat\alpha(y) \neq \hat\alpha(x)
\]

\page{31}
On dit alors que le jeu \struck{de} de position à deux joueurs est un jeu
\emph{alterné} (les deux joueurs jouent alternativement). Pour qu'une
application $\hat\alpha$ existe prolongeant $\alpha$ et ayant les propriétés
précédentes, il faut et suffit que les deux conditions suivantes soient
satisfaites
\begin{enumerate}
  \item $\forall x, y \in C$, tels que $y \in R(x)$, $y \notin C_0$, on a
  $\alpha(y) \neq \alpha(x)$
  \item $\forall y \in C_0$, soit
  $T(y) = \{ x \in C \mid y \in R(x) \} \subset C \setminus C_0$, alors
  $\alpha \mid T(y)$ est constante.
\end{enumerate}

Sous les conditions de 2), on aura donc nécessairement, si
$T(y) \neq \emptyset$ et $\alpha(T(y)) = \{ j \}$, $\hat\alpha(y) = j'$, où
$j'$ est l'élément de $J$ qui est $\neq j$. Cela montre que $\hat\alpha$ est
bien déterminé sur l'ens.\ des $y \in C_0$ tels que $T(y) \neq \emptyset$,
donc aussi sur l'ens.\ des positions \emph{non isolées} (une position $y$ est
dite \emph{isolée} si elle n'a pas de position antécédente ni suivante, i.e.\
s'il n'existe pas de $x \in C$ tel qu'on ait $y \in R(x)$ ou $x \in R(y)$). On
peut attribuer à $\hat\alpha$ des valeurs arbitraires sur les positions
isolées.

Quand on a défini un $\hat\alpha$, on peut définir
\[
\begin{array}{l}
G = \{ x \in C \mid x \in G(\hat\alpha(x)) \} \\
P = \{ x \in C \mid x \in G(\bar{\hat\alpha}(x)) \}
\end{array}
\]
et on \struck{\ill{}} aura

\page{32}
\[
\begin{cases}
G \cap (C \setminus C_0) = G^{*} \\
P \cap (C \setminus C_0) = P^{*} \\
C \setminus (G \cup P) = N
\end{cases}
\]
(positions nulles, i.e.\ gagnantes pour aucun des deux joueurs)

Les positions $x$ seront encore dites \emph{gagnantes} (resp.\
\emph{perdantes}) si $x \in G$ (resp.\ $x \in P$) — cette terminologie
coïncide avec la précédente si $x \notin C_0$. On aura, posant
\struck{\ill{}}
\[
\hat C(j) = \{ x \in C \mid \hat\alpha(x) = j \}
\]
(de sorte que $C$ apparaît comme réunion disjointe des deux ens.\
$\hat C(j)$, $j \in J$),
\[
\begin{cases}
G(j) \cap \hat C(j) = \hat C(j) \cap G \\
G(j) \cap \hat C(\bar\jmath) = \hat C(\bar\jmath) \cap P
\end{cases}
\]
donc les $G(j)$ ($j \in J$) se reconstituent à partir de $\hat\alpha$ et des
ensembles $G$, $P$.

On aura, pour $x \in C \setminus C_0$
\[
\begin{array}{lll}
x \in G & \Longleftrightarrow \exists\, y \in R(x) \text{ tel que } y \in P & \text{i.e. } R(x) \cap P \neq \emptyset \\
x \in P & \Longleftrightarrow \forall\, y \in R(x), \text{ on a } y \in G & \text{i.e. } R(x) \subset G
\end{array}
\]
i.e.
\[
\begin{array}{l}
G \cap (C \setminus C_0) = \{ x \in C \setminus C_0 \mid R(x) \cap P \neq \emptyset \} \\
P \cap (C \setminus C_0) = \{ x \in C \setminus C_0 \mid R(x) \subset G \}
\end{array}
\]

\page{33}
\emph{Proposition} Soit
$(C, R, J, \hat\alpha, G_0 : J \to \mathfrak{P}(C_0))$ un jeu de positions à
deux joueurs jouant alternativement, $G \subset C$ et $P \subset C$, l'ens.\
des positions gagnantes resp.\ perdantes. Soient $\widetilde G$,
$\widetilde P$ deux parties de $C$ ayant les propriétés suivantes
\begin{enumerate}
  \item $\forall x \in \widetilde G$, $\exists\, y \in R(x)$ tel que
  $y \in \widetilde P$
  \item $\forall x \in \widetilde P$, $\forall y \in R(x)$, on a
  $y \in \widetilde G$
  \item $\forall x \in \widetilde G \cap C_0$, on a
  $x \in G_0(\hat\alpha(x))$
  \item $\forall x \in \widetilde P \cap C_0$, on a
  $x \in G_0(\bar{\hat\alpha}(x))$
\end{enumerate}
Alors, si dans $C$ il n'existe pas de \struck{\ill{}} partie infinie, on a
\[
\widetilde G \subset G, \qquad \widetilde P \subset P
\]
Si de plus on suppose \struck{\ill{} disjointes des deux parties $G_0(j)$
($j \in J$)} que $\widetilde G$ et $\widetilde P$ sont complémentaires l'un
de l'autre, \struck{\ill{}} et $G_0(j) \cap G_0(\bar\jmath) = \emptyset$
(donc $G \cap P = \emptyset$) alors on a $\widetilde G = G$,
$\widetilde P = P$ (et par suite $G \cup P = C$ donc
$G_0(j) \cup G_0(j') = C_0$)

\emph{Cor} Soit $(C, \dots)$ un jeu alterné tel que \struck{\ill{}} a) Pas de
partie infinie b) $C_0$ réunion disjointe de $G_0(j)$ et $G_0(\bar\jmath)$.
Soient $\widetilde G$, $\widetilde P$ deux parties de $C$. Pour que
$\widetilde G = G$, $\widetilde P = P$, il faut et il suffit que
$\widetilde G$ et $\widetilde P$ soient complémentaires l'une de l'autre, et
satisfont aux conditions 1) à 4) ci-dessus.

\section{Exponentielles (en car.~0)}
\note{le titre est de sa main, en tête de la page 35.}

\page{35}
\subsection*{I) Généralités}

$\mathbb{A}$ $\mathbb{Q}$-algèbre (ass.\ avec 1)

$N(\mathbb{A})$ ens.\ des él.\ nilp.\ de $\mathbb{A}$. On a donc
$1 + N(\mathbb{A}) \subset \mathbb{A}^{*}$ (ens.\ des él.\ inv.\ de
$\mathbb{A}$)
\[
N(\mathbb{A}) \xrightarrow{\;\exp\;} 1 + N(\mathbb{A}) \qquad \exp x = \sum_{n \geq 0} \frac{x^n}{n!}
\]

\emph{Th 1} 1) $\exp$ est bijective. L'application inverse est donnée par
\[
\log(1+y) = \sum_{n \geq 1} (-1)^{n+1} \frac{y^n}{n} \qquad (y \in N(\mathbb{A}))
\]
2) Si $x, y \in N(\mathbb{A})$ commutent, alors $x + y \in N(\mathbb{A})$ et
\[
\exp(x+y) = \exp x \, \exp y \qquad \text{NB } \exp 0 = 1
\]
Si $u, v \in 1 + N(\mathbb{A})$ commutent, alors $uv \in 1 + N(\mathbb{A})$,
et
\[
\log(uv) = \log u + \log v \qquad \text{NB } \exp 1 = 0
\]
\note{sic : « $\exp 1 = 0$ » ; on attend $\log 1 = 0$.}

3) Soit $L \subset \mathbb{A}$ un sous-$\mathbb{Q}$-espace vectoriel tel que
$x \in L$, $n \in \mathbb{N}^{*}$ implique $x^n \in L$, \struck{Alors} et $L$
formé d'éléments nilpotents. Alors
\[
x \mapsto \exp x : L \xrightarrow{\;\sim\;} 1 + L
\]

\emph{Remarques} Soit $L \subset \mathbb{A}$ un sous-$\mathbb{Q}$-espace
vectoriel de $A$ \struck{les conditions de 3), supposons que} tel que
$L \subset N(A)$, et tel que $x, y \in L \Rightarrow xy \in L$. Alors
$L \xrightarrow{\sim} 1 + L$ par 3°, et de plus $1 + L$ est un sous-groupe de
$\mathbb{A}^{*}$. D'où par transport de str.\ une str.\ de groupe sur $L$.
Elle est commutative \struck{si} ssi $L$ est commutatif, i.e.\ si
$x, y \in L \Rightarrow xy = yx$ — et dans ce cas cette structure est
l'\emph{addition} de $L$.
\marginal{réciproque ? \ill{}}

\emph{Exemple 1} Soient $K$ une $\mathbb{Q}$-alg.\ commut.\ \struck{\ill{}},
$\mathbb{A}$ une $K$-algèbre (ass.\ avec 1), $J$ un idéal nil de $K$. Alors
$J\mathbb{A}$ satisfait les conditions de la Remarque, et on a donc
\[
J\mathbb{A} \longrightarrow 1 + J\mathbb{A}
\]
\marginal{$\alpha : \mathbb{A} \to \mathbb{A}_0$ \uncertain{can.} Alors
\struck{$J\mathbb{A} =$}}

\emph{NB} Soient $K_0 = K/J$,
$\mathbb{A}_0 = \mathbb{A}/J\mathbb{A} = \mathbb{A} \otimes_K K_0$,
\[
J\mathbb{A} = \{ x \in \mathbb{A} \mid \alpha(x) = 0 \}, \qquad
1 + J\mathbb{A} = \{ x \in \mathbb{A} \mid \alpha(x) = 1 \}
\]
mais avec $J$ \emph{nilpotent}, \ill{}

\emph{Exemple 2} $K$, $J$ comme dans Exemple 1, $M$ un $K$-module,
$M_0 = M \otimes_K K_0 = M/JM$, $\alpha_M : M \to M_0$ appl.\ can. ;
$\mathcal{U} = \operatorname{End}_K(M)$, \struck{\ill{}} posons, alors
\[
L = \{ u \in \mathcal{U} = \operatorname{End}_K(M) \mid u \otimes_K K_0 = 0 \} = \{ u \mid u(M) \subset JM \}
\]
\[
1 + L = \{ u \in \mathcal{U} = \operatorname{End}_K(M) \mid u \otimes_K K_0 = \mathrm{id}_{M_0} \}.
\]
\struck{Ceci posé}
Ceci posé, les conditions de la remarque sont satisfaites donc on a
\[
\exp : L \xrightarrow{\;\sim\;} 1 + L
\]

\page{36}
\subsection*{II) Groupes à un paramètre formels}

$K$ $\mathbb{Q}$-algèbre comm.\ unifère.

$\overline{\mathbb{G}}_{a,K}$ droite formelle sur $K$
($\overline{\mathbb{G}}_{a,K} \subset \mathbb{G}_{a,K}$) :
\[
\overline{\mathbb{G}}_{a}(K') = N(K') \subset K' \qquad (K' \ K\text{-alg.\ comm.} \dots)
\]

1) Soit $\mathcal{U}$ $K$-algèbre, $\mathrm{Un}(\mathcal{U})$ le
foncteur-groupes
\[
K' \mapsto (\mathcal{U}_{K'})^{*} \qquad \mathrm{Alg}_{/K} \to \mathrm{Gr}
\]
On va \uncertain{définir} les homomorphismes de groupes
\[
\begin{array}{l}
\overline{\mathbb{G}}_{a,K} \to \mathrm{Un}(\mathcal{U}) \\
\mathbb{G}_{a,K} \to \mathrm{Un}(\mathcal{U})
\end{array}
\]
Soit $x \in \mathcal{U}$, on lui associe
\[
\bar e_x : \overline{\mathbb{G}}_{a,K} \to \mathrm{Un}(\mathcal{U})
\]
par
\[
\bar e_x(\lambda) = \exp(\lambda x)
\]
i.e.\ si $\lambda \in \overline{\mathbb{G}}_{a,K}(K') = N(K')$,
$\bar e_x(\lambda) = \exp(\lambda x_{K'})$ (cf.\ Ex.~1). On a bien
$\bar e_x(\lambda + \lambda') = \bar e_x(\lambda)\, \bar e_x(\lambda')$
(Th.~1, 2°). Donc
\[
\bar e_x \in \mathrm{Hom}_{\mathrm{gr}}(\overline{\mathbb{G}}_{a,K}, \mathrm{Un}(\mathcal{U}))
\]
D'où une application
\[
(*) \qquad x \mapsto \bar e_x : \mathcal{U} \to \mathrm{Hom}_{\mathrm{gr}}(\overline{\mathbb{G}}_{a,K}, \mathrm{Un}(\mathcal{U}))
\]
Si $x \in N(\mathcal{U})$, alors on définit de même
\[
\begin{cases}
e_x : \mathbb{G}_{a,K} \to \mathrm{Un}(\mathcal{U}) \\
e_x(\lambda) = \exp(\lambda x)
\end{cases}
\]
et on a encore
\[
e_x \in \mathrm{Hom}_{\mathrm{gr}}(\mathbb{G}_{a,K}, \mathrm{Un}(\mathcal{U}))
\]
d'où
\[
(**) \qquad x \mapsto e_x : N(\mathcal{U}) \to \mathrm{Hom}_{\mathrm{gr}}(\mathbb{G}_{a,K}, \mathrm{Un}(\mathcal{U}))
\]

\emph{NB} a) $\bar e_x$ est la restriction de $e_x$ à
$\overline{\mathbb{G}}_{a,K}$. b) Si $x, y$ commutent, alors $\bar e_x$,
$\bar e_y$ (resp.\ $e_x$, $e_y$) commutent, et
$\bar e_{x+y} = \bar e_x \cdot \bar e_y$ (resp.\
$e_{x+y} = e_x \cdot e_y$).

\emph{Th.~2} $(*)$ et $(**)$ sont \emph{bijectives}. [NB si $\mathcal{U}$ est
commutatif alors ce sont des hom.\ de groupes]

\page{37}
\[
e_x(\lambda) = \sum \frac{x^n}{n!} \lambda^n
\]
d'où
\[
e_x = e_y \Longleftrightarrow \frac{x^n}{n!} = \frac{y^n}{n!} \ \ \forall n \in \mathbb{N} \Longleftrightarrow x = y
\]
donc $x \mapsto e_x$ injectif
$N(\mathcal{U}) \to \mathrm{Hom}_{\mathrm{gr}}(\mathbb{G}_{a,K}, \mathrm{Un}(\mathcal{U}))$.
On voit de même que $x \mapsto \bar e_x$
$\mathcal{U} \to \mathrm{Hom}_{\mathrm{gr}}(\overline{\mathbb{G}}_{a,K}, \mathrm{Un}(\mathcal{U}))$
est injectif. (\ill{} \uncertain{mult.})

Un hom.\ quelconque
$\bar\varphi : \overline{\mathbb{G}}_a \to \mathrm{Un}(\mathcal{U}) \subset$
\uncertain{$W(\mathcal{U})$} est donné par
\[
\bar\varphi(\lambda) = \sum_{n \geq 0} a_n \lambda^n \qquad a_n \in \mathcal{U}
\]
La condition $\bar\varphi(0) = 1$ signifie $a_0 = 1$. Écrivons
$\bar\varphi(\lambda + \mu) = \bar\varphi(\lambda)\, \bar\varphi(\mu)$ i.e.
\[
\begin{array}{ccc}
\sum a_n (\lambda + \mu)^n & = & \sum a_i \lambda^i \sum a_j \mu^j \\
\| & & \| \\
\sum_{i,j} a_{i+j} \frac{(i+j)!}{i!\,j!} \lambda^i \mu^j & & \sum a_i a_j \lambda^i \mu^j
\end{array}
\]
on trouve les conditions $\forall n \in \mathbb{N}$, $i, j \in \mathbb{N}$
tels que $i + j = n$
\[
\frac{(i+j)!}{i!\,j!}\, a_n = a_i a_j
\]
\note{au-dessus, une première écriture barrée : $a_n (= a_i a_j)$.}
Donc, faisant $i = 1$, $j = n - 1$
\[
n a_n = a_1 a_{n-1} \quad \text{i.e.} \quad a_n = \tfrac{1}{n} a_1 (a_{n-1})
\]
d'où par récurrence sur $n$
\[
a_n = \frac{1}{n!} x^n \qquad x = a_1
\]
(Condition aussi suffisante …). Donc $\bar\varphi = \bar e_x$.

Soit $\varphi : \mathbb{G}_{a,K} \to \mathrm{Un}(\mathcal{U}) \subset W(\mathcal{U})$,
\struck{\ill{}} on a encore
\[
\varphi(\lambda) = \sum a_n \lambda^n \qquad a_n \in \mathcal{U}, \text{ les } a_n \text{ nuls sauf un nb fini}
\]
Utilisant $\bar\varphi = \varphi \mid \overline{\mathbb{G}}_a$, on trouve
$a_n = \frac{x^n}{n!}$. Comme $a_n = 0$ pour $n$ grand, on trouve $x$
nilpotent, cqfd.

\page{38}
2) Soit $M$ un $K$-module, $\mathbf{Aut}_K(M)$
$[\subset \mathbf{End}_K(M)]$ le foncteur en groupes
$K' \mapsto \mathrm{Aut}_{K'}(M_{K'})$ sur $\mathrm{Alg}_{/K}$. On cherche
les hom.\ de groupes
\[
\begin{array}{l}
\mathbb{G}_a \xrightarrow{\;\varphi\;} \mathbf{Aut}_K(M) \\
\overline{\mathbb{G}}_a \xrightarrow{\;\bar\varphi\;} \mathbf{Aut}_K(M)
\end{array}
\]
On trouve comme dans 1)
\[
\begin{array}{ll}
(*) & u \mapsto \bar e_u : \mathrm{End}_K(M) \to \mathrm{Hom}_{\mathrm{gr}}(\overline{\mathbb{G}}_{a,K}, \mathbf{Aut}_K(M)) \\
(**) & u \mapsto e_u : N(\mathrm{End}_K(M)) \to \mathrm{Hom}_{\mathrm{gr}}(\mathbb{G}_{a,K}, \mathbf{Aut}_K(M))
\end{array}
\]
\[
\begin{array}{ll}
\bar e_u(\lambda) = \exp(\lambda u) = \sum_{n \geq 0} \frac{u^n_{K'}}{n!} \lambda^n & \lambda \in \overline{\mathbb{G}}_a(K') = N(K') \\
e_u(\lambda) = \text{— — —} & (\lambda \in \mathbb{G}_a(K') = K')
\end{array}
\]
On trouve encore que $(*)$ et $(**)$ sont des bij.

(\emph{NB} Si $M$ est projectif de type fini, il suffirait d'appliquer le cas
précédent à $\mathcal{U} = \mathrm{End}_K(M)$)

On utilise le fait que les morphismes de foncteurs
$\bar\varphi : \overline{\mathbb{G}}_{a,K} \to \mathbf{End}_K(M)$ (pas nec.\
multiplicatif) est donné de façon unique par
\[
\bar\varphi(\lambda) = \sum_{n \geq 0} u_n \lambda^n \qquad u_n \in \mathrm{End}_K(M),
\]
et les morphismes de foncteurs
$\varphi : \mathbb{G}_{a,K} \to \mathbf{End}_K(M)$ par
\[
\varphi(\lambda) = \sum_{n \geq 0} u_n \lambda^n \qquad u_n \in \mathrm{End}_K(M) \text{ nuls sauf un nb fini.}
\]
Ceci posé, \struck{le calcul} on constate que la condition néc.\ pour que
$\bar\varphi$ (resp.\ $\varphi$) soit un hom.\ de groupes, c'est qu'on ait
$u_n = \frac{1}{n!}(u_1)^n$ …

\page{39}
\emph{Prop} Soient $M, N$ deux $K$-modules, $u$ ($v$) un endom.\ de $M$
($N$), $f : M \to N$ un $K$-hom. Considérons $M, N$ comme
$\overline{\mathbb{G}}_{a,K}$-modules par
$\lambda \in \overline{\mathbb{G}}_{a,K}(K')$ opérant par $\exp(\lambda u)$
resp.\ $\exp(\lambda v)$. Pour que $f$ soit un
$\overline{\mathbb{G}}_{a,K}$-morphisme, il faut et suffit que l'on ait
$vf = fu$ i.e.\ $vf - fu = 0$.

\emph{Cor} Si $u, v$ sont nilpotents, c'est aussi la cond.\ néc.\ et suff.\
pour que $f$ soit un $\mathbb{G}_a$-morphisme.

\emph{Dém} Il faut écrire
\[
\sum_{n \geq 0} \lambda^n \frac{v^n}{n!} f = \sum_{n \geq 0} \lambda^n f \frac{u^n}{n!}
\]
\uncertain{identiquement} en $\lambda$ nilp.\ i.e.
\[
\sum_{n \geq 1} \frac{1}{n!} (v^n f - f u^n) \lambda^n = 0
\]
\uncertain{identiquement} en $\lambda$ nilp., ce qui équivaut à :
\[
v^n f - f u^n = 0 \qquad \forall n \geq 1
\]
Pour $n = 1$, c'est la commutation $vf = fu$, qui implique par récurrence
$v^n f = f u^n$ ($\forall n \geq 1$).

\emph{Proposition} Soient $u \in N(\mathrm{End}_K(M))$,
$v \in N(\mathrm{End}_K(N))$, $f \in \mathrm{Hom}_K(M, N)$. Conditions
équivalentes :
\begin{enumerate}
  \item $f$ est un $\overline{\mathbb{G}}_{a,K}$-hom.\ (pour
  $\overline{\mathbb{G}}_a$ opérant sur $M, N$ par $\bar e_u$ et
  $\bar e_v$)
  \item $f$ est un $\mathbb{G}_{a,K}$-hom.\ (…)
  \item \struck{$fu$} $= vf$ \struck{\ill{}} i.e.\
  $\delta \overset{\text{déf}}{=} vf - fu = 0$
  \item $f \circ \exp(u) = \exp(v) \circ f$
\end{enumerate}

\page{40}
\emph{Dém} On a déjà vu l'équivalence de (1) à (3), et
\struck{\ill{}} cela implique 4 (i.e.\ $e_v(1) f = f e_u(1)$). Prouvons
$(4) \to (3)$. Or (4) s'écrit
\[
\underbrace{(\exp v) \circ f - f \circ \exp u}_{\displaystyle \sum_{n \geq 1} \frac{1}{n!} \underbrace{(v^n f - f u^n)}_{\delta_n \ (\text{déf})}} = 0
\]
Or on a (réc.\ sur $n$)
\[
\delta_n = v^{n-1} \delta + v^{n-2} \delta u + \dots + v \delta u^{n-2} + \delta u^{n-1}
\]
donc (4) signifie
\[
\delta + \frac{1}{2}(v \delta + \delta u) + \frac{1}{3!}(v^2 \delta + v \delta u + \delta u^2) + \dots = 0 \qquad (\text{somme finie})
\]
On peut l'écrire
\[
\delta = \sum_{i+j \geq 1} c_{ij}\, v^i \delta u^j \qquad \Bigl(c_{ij} \in K,\ c_{ij} = \frac{1}{(i+j)!}\Bigr)
\]
\note{tel qu'écrit ; l'équation qui précède donne
$c_{ij} = -\frac{1}{(i+j+1)!}$. Seule compte ici l'existence des
$c_{ij} \in K$.}
Or soit
\[
\mathcal{U} = \mathrm{End}_K(M)
\]
\[
\Theta, \Theta' : \mathcal{U} \to \mathcal{U} \qquad K\text{-endom.}
\]
\[
\Theta(x) = vx, \qquad \Theta'(x) = xu
\]
Donc $\Theta$, $\Theta'$ sont nilpotents ($v$, $u$ l'étant) et commutent, et
on a
\[
\delta = \sum_{i+j \geq 1} c_{ij}\, \Theta^i \Theta'^j \delta = \rho \delta
\]
\[
\rho = \sum_{i+j \geq 1} c_{ij}\, \Theta^i \Theta'^j \in \mathrm{End}_K(\mathcal{U})
\]
où $\rho$ est \emph{nilpotent} ($\Theta$, $\Theta'$ l'étant et commutant). Or
$\delta = \rho \delta$ implique par récurrence
\note{la phrase se poursuit au lot suivant.}

\end{document}
